\section{Introduction}
\label{sec:introduction}

Every US state publishes early learning standards: the developmental
expectations that shape curriculum, assessment, and teacher preparation for
children from birth through kindergarten entry. More than fifty jurisdictions
issue them independently~\citep{ncecqa_elgs,ecs2024governance}, and each designs its
own document. The result is a body of policy that is public, consequential, and
almost entirely unavailable to software. Unlike K--12 academic standards, which
have machine-readable distribution
formats~\citep{sutton2008achievement,1edtech2017case}, early learning standards
exist as PDFs --- multi-column tables, age-banded grids, and prose narratives
whose typography is shared by no two states.

The obstacle is not optical character recognition. It is that a standards
document is a \emph{tree} rendered as a page. A single indicator is meaningful
only through its ancestry: which domain, which strand, which sub-strand it sits
under, and at what age band. Recovering that tree requires deciding, for every
line, how deeply it is nested --- and the surface cues that would answer the
question are exactly what varies between states. Arizona writes
\texttt{Strand 1: Language and Literacy}; Kentucky writes
\texttt{Benchmark 1.1}; Nevada prints a code namespace that skips a level
entirely; Colorado publishes no sub-strand tier at all. A rule tuned to one
document's vocabulary is, by construction, a rule that fails on the next.

My central claim is that a language model can recover this structure by
classifying each element's \textbf{nesting position} rather than its
document-specific label, and that doing so is what generalizes. Concretely, I
first infer a per-document \emph{depth map} --- how many nesting levels this
particular document uses, and what each looks like --- and then classify every
element against that map as an authority, rather than against a fixed
vocabulary of level names. The label \texttt{Benchmark 1.1} then does not have
to be recognized as a sub-strand; it has to be recognized as sitting one level
below whatever \texttt{Language and Early Literacy Standard 1} is.

This framing also determines how the system is built. I hold a deliberate
design discipline that structural decisions belong in the prompt as general
document principles, not in Python as per-document branches. That discipline
was adopted after an earlier rule-driven revision scored well on the states it
had been developed against and failed to transfer; the migration that removed
those rules is described in \S\ref{sec:method-llm-first}. Deterministic code is
retained only where it reads the \emph{shape} of an output --- dot-segment
structure, whitespace, string containment --- and never any document's words.

The absence of a labeled corpus is what rules out the obvious supervised
alternative. There is no annotated dataset of hierarchically-parsed early
learning standards at any scale, so a fine-tuned model would first require one
to be built, and the dominant cost of that path is annotation rather than
compute. My own goldens took expert effort to produce and cover six states;
they are an evaluation instrument, not a training set. An in-context approach
that needs no labeled data is therefore not merely cheaper here, it is the only
approach available without first solving a different problem
(cf.~\citealp{liu2022fewshot} for the counterargument where labels do exist).

\paragraph{Contributions.}
\begin{enumerate}
  \item A two-stage detection method built on per-document depth-map inference
    and classify-by-nesting-position prompting, with no per-state code path
    (\S\ref{sec:method-detection}; what the prompts do carry is itemized in
    \S\ref{sec:method-discipline}).
  \item A parsing stage that resolves detected elements into a single canonical
    schema across documents that do not agree on level count or code namespace,
    producing deterministic, collision-free identifiers
    (\S\ref{sec:method-parsing}).
  \item A serverless architecture that batches detection and parsing through a
    prepare--map--merge pattern so the method applies to whole standards
    documents rather than excerpts (\S\ref{sec:method-deployment},
    Appendix~\ref{sec:appendix-architecture}).
  \item An evaluation over four development states and two held-out states.
    Detector recall is 1.000 at every hierarchy level in all six states, and a
    manual audit of all 297 in-scope detections finds a single hallucinated
    element (verified precision 0.997). An ablation shows the depth-map pass is
    what carries the intermediate levels: removing it leaves domain and
    indicator recall untouched while strand falls to 0.920 and sub-strand to
    0.833 (\S\ref{sec:experiments-ablation}).
  \item A reporting discipline that the rest of the evaluation is built to
    satisfy, and which I offer as a contribution rather than as a caveat: every
    table is labelled with the corpus tier it was computed at; every unmatched
    detection is audited by hand and signed rather than assumed to be a false
    positive; every number is regenerated from a recorded artifact by a named
    script; a validation guard is shown catching a live defect rather than
    asserted to work; and a previously recorded result of my own that failed to
    reproduce is reported as a non-reproduction instead of quietly replaced
    (\S\ref{sec:experiments-stability}, \S\ref{sec:discussion-limitations}).
\end{enumerate}

\paragraph{Scope of the reported numbers.} Every headline quality measurement
in this paper is computed at the \emph{subset} corpus tier: hand-trimmed
excerpts of one to two domains per state, 70 pages in total, fully annotated.
This is favorable preprocessing and I state it plainly rather than leaving it to
be inferred; \S\ref{sec:corpus} gives the tiers and retained page ranges, and
\S\ref{sec:discussion-limitations} discusses what does and does not transfer to
whole documents. Cost, scale and one quality measurement are reported at the
\emph{trimmed} tier instead---a document's standards content in full, with front
matter, expository essays and appendices removed. The tiers are never mixed
within a table, and each table names its own
(\S\ref{sec:experiments-scale}, \S\ref{sec:experiments-scale-quality}).
